\documentclass[sigplan,10pt]{acmart}

\setcopyright{none}
\acmYear{2026}
\acmConference[ATC '26]{2026 ACM Annual Technical Conference}{2026}{}
\acmBooktitle{2026 ACM Annual Technical Conference}
\acmDOI{}
\acmISBN{}
\settopmatter{printfolios=true,printacmref=false}
\renewcommand\footnotetextcopyrightpermission[1]{}
\pagestyle{plain}

\usepackage{booktabs}

\newcommand{\system}{GPUTriage}

\begin{document}
\sloppy

\title{LLM-Based Triage of GPU Numerical Failures from Public Issue Reports}

\author{Paper \#XXX}
\affiliation{\institution{Anonymous Submission}\country{Anonymous}}

\maketitle

\section{Motivation}

GPU software stacks now combine tensor libraries, generated kernels, compiler
passes, mixed-precision execution, and hardware-specific runtime behavior.
When these systems fail, the first useful evidence often appears in public
issue trackers as noisy natural-language reports: NaNs, Inf values, overflow,
wrong answers, dtype conversion failures, compiler crashes, and performance
regressions are frequently mixed in the same thread. This creates a practical
systems problem. Maintainers need to route reports to the right experts, and
researchers need reusable traces of reliability failures, but weak labels from
keyword search are not trustworthy enough to support either goal.

\section{Contribution}

We present \system{}, a human-labeled benchmark and LLM-based triage system for
GPU numerical-failure reports mined from public repositories. The benchmark
contains 1,191 human-labeled issues from eight GPU-related projects spanning
compiler, runtime, tensor, and data-processing software. Each issue receives one
primary label from a seven-class taxonomy covering invalid values,
overflow/underflow, precision or tolerance failures, dtype or casting failures,
crash/compile failures, performance-only reports, and non-numerical reports.
The artifact also includes annotation instructions, blind-audit files,
adjudication logs, taxonomy repair logs, model outputs, evaluation tables, an
error appendix, and root-cause evidence files.

\begin{table}[t]
\centering
\caption{Key ATC submission results.}
\label{tab:ea_results}
\begin{tabular}{lcc}
\toprule
Setting & Accuracy & Macro F1 \\
\midrule
Weak retrieval labels & 0.505 & 0.426 \\
Voting ensemble & 0.738 & 0.712 \\
FLAN-T5-base LLM & 0.805 & 0.778 \\
LLM/ensemble agreement & 0.893 & 0.844 \\
Leave-one-repository-out & 0.668 & 0.648 \\
\bottomrule
\end{tabular}
\end{table}

\section{Method and Results}

\system{} separates data mining from ground truth. Candidate issues are first
retrieved with terms such as \texttt{nan}, \texttt{inf}, \texttt{overflow},
\texttt{dtype}, \texttt{precision}, and \texttt{wrong result}. These candidate
labels are treated only as weak supervision. Human annotation then repairs the
label space and creates the canonical benchmark. A blind audit over 119 rows
achieves 80.7\% agreement and Cohen's $\kappa=0.721$; subsequent adjudication
passes apply 419 human updates and change 211 primary labels.

This design targets an operational systems workflow rather than a generic text
classification benchmark. A maintainer-facing assistant should identify which
reports are likely numerical failures, route them into the right queue, and
abstain when the evidence is ambiguous. For that reason, we evaluate both
full-coverage prediction and selective prediction. The latter is important
because wrong labels are costly: a dtype report routed as a compiler crash or a
non-numerical API question routed as a correctness bug wastes maintainer time.

The system evaluates both deterministic models and a fine-tuned LLM. Weak
retrieval labels reach only 50.5\% accuracy, demonstrating that raw
issue-search labels are not reliable ground truth. A deterministic voting
ensemble reaches 73.8\% accuracy and 0.712 macro F1. A fine-tuned FLAN-T5-base
classifier reaches 80.5\% held-out accuracy and 0.778 macro F1. For
high-confidence triage, an LLM/deterministic agreement mode reaches 89.3\%
accuracy and 0.844 macro F1 at 68.3\% coverage. Chronological evaluation
reaches 76.6\% accuracy, while leave-one-repository-out evaluation reaches
66.8\%, exposing cross-project transfer as the main remaining challenge.

\section{Why This Fits ATC}

The work is a systems paper because it studies practical failure reports from
real GPU software ecosystems and produces an artifact for bug finding,
analyzing, and troubleshooting. Its main result is not only that an LLM can
classify issue text; it is that reliable triage requires curated systems data,
auditable label repair, cross-repository evaluation, and selective confidence
mechanisms. The benchmark makes previously unavailable GPU numerical-failure
traces available for future systems reliability work. The full paper discusses
the system design, data construction, model evaluation, cross-repository
weaknesses, and limitations needed to judge whether the results support the
claim.

The expected ATC contribution is therefore an artifact-backed systems result:
public GPU issue reports can be turned into a reusable triage benchmark, but
only after weak labels are repaired and cross-repository transfer is measured
explicitly. The paper does not claim definitive root-cause diagnosis from issue
text alone. Instead, it establishes report-level triage as a measurable systems
task and provides the data needed for future root-cause and debugging research.

\bibliographystyle{ACM-Reference-Format}
\begin{thebibliography}{4}

\bibitem{triton2019}
Tillet, P., Kung, H., and Cox, D. 2019. Triton: An intermediate language and
compiler for tiled neural network computations. In \emph{MAPL}.

\bibitem{pytorch2019}
Paszke, A. et al. 2019. PyTorch: An imperative style, high-performance deep
learning library. In \emph{NeurIPS}.

\bibitem{tvm2018}
Chen, T. et al. 2018. TVM: An automated end-to-end optimizing compiler for deep
learning. In \emph{OSDI}.

\bibitem{t5}
Raffel, C. et al. 2020. Exploring the limits of transfer learning with a
unified text-to-text transformer. \emph{JMLR} 21, 140, 1--67.

\end{thebibliography}

\end{document}
